\section*{Bab \ref{ch:b2:atmospheric-thermodynamics} --- Termodinamika Atmosfer}

\begin{solution}{exo:b2:atmospheric-thermodynamics:1}
(a) $e = e_s(\qty{278.15}{K}) = \qty{8.73}{hPa}$ dan $e_s(\qty{283.15}{K}) = \qty{12.32}{hPa}$.
(b) $8.73/12.32 = \qty{70.9}{\percent}$. (c) $r = 0.622 \times 8.73/(850 - 8.73) =
\qty{6.45}{g/kg}$ dan $q = r/(1 + r) = \qty{6.41}{g/kg}$.
\end{solution}

\begin{solution}{exo:b2:atmospheric-thermodynamics:2}
(a) $\mathrm{d}\ln e_s/\mathrm{d}T = 5419/T^2$: \qty{7.26}{\percent} per kelvin pada
\qty{273.15}{K}, \qty{6.31}{\percent} pada \qty{293.15}{K}, \qty{5.90}{\percent} pada
\qty{303.15}{K}. (b) $e_s(\qty{20}{\celsius}) = \qty{23.66}{hPa}$ dan $e_s(\qty{25}{\celsius})
= \qty{32.26}{hPa}$; rasionya 1.363, dibandingkan dengan $1.063^5 = 1.357$. Laju itu adalah turunan
$\ln e_s$, sehingga selama \qty{5}{K} faktornya $\exp(5 \times 0.0620) = 1.363$, dengan 0.0620 sebagai
laju pada titik tengah \qty{22.5}{\celsius}; menggandakan \qty{6.3}{\percent} lima kali secara majemuk memberikan nilai yang dekat
tetapi tidak eksak. (c) $23.66/23.39 = 1.012$: rumus itu \qty{1.2}{\percent} terlalu tinggi, karena
rumus itu mempertahankan kalor laten titik tripel, sedangkan kalor laten berkurang ketika suhu
naik; kemiringan rumus itu terlalu curam, dan kesalahannya bertambah dengan jarak dari
\qty{0}{\celsius}.
\end{solution}

\begin{solution}{exo:b2:atmospheric-thermodynamics:3}
(a) $125 \times (30 - 14) = \qty{2000}{m}$. (b) $30 - 9.8 \times 2.0 = \qty{10.4}{\celsius}$.
(c) $14 - 1.8 \times 2.0 = \qty{10.4}{\celsius}$: suhu dan titik embun bertemu, sebagaimana seharusnya pada
aras kejenuhan; aturan itu adalah perhitungan ini. (d) $125 \times (33 - 14) = \qty{2375}{m}$:
dasar awan naik ketika sore hari menghangatkan udara dekat tanah.
\end{solution}

\begin{solution}{exo:b2:atmospheric-thermodynamics:4}
(a) Kurva titik embun berimpit dengan kurva suhu dari tanah sampai sekitar \qty{908}{hPa}:
setengah kilometer terbawah jenuh, atau hampir jenuh. (b) Antara sekitar 908 dan \qty{889}{hPa}
kurva suhu membelok ke kanan, ke nilai yang lebih hangat, ketika tekanan berkurang: suhu
naik terhadap ketinggian. (c) LCL pada sekitar \qty{964}{hPa}, beberapa puluh meter di atas tanah; LFC pada
sekitar \qty{730}{hPa}; EL pada sekitar \qty{235}{hPa}. (d) Di puncak inversi titik embun
\qty{11}{K} di bawah suhu, dan \qty{20}{K} di bawahnya pada \qty{850}{hPa}: udara hangat yang sangat kering
terletak di atas udara sejuk yang lembap. Kelembapan relatifnya rendah, sehingga udara itu tidak memperoleh
kehangatannya dari kondensasi; udara yang turun menghangat karena pemampatan dengan laju adiabatik kering sambil
mempertahankan \omterm{def:b2:atmospheric-thermodynamics:moist-air}{rasio campurannya}, dan hal itu akan menghasilkan lapisan seperti itu.
\end{solution}

\begin{solution}{exo:b2:atmospheric-thermodynamics:5}
(a) $T_v = 303.15 \times (1 + 0.020/0.622)/1.020 = \qty{306.76}{K}$; hampiran memberikan
$303.15 \times (1 + 0.61 \times 0.020) = \qty{306.85}{K}$. (b) $\rho = p/(R_dT_v) = 10^5/(287.0
\times 306.76) = \qty{1.1358}{kg/m^3}$, dibandingkan dengan $10^5/(287.0 \times 303.15) =
\qty{1.1494}{kg/m^3}$ untuk udara kering: \qty{1.2}{\percent} lebih ringan. (c) Sebesar $T_v - T =
\qty{3.6}{K}$. (d) Parsel yang naik hanya perlu lebih ringan daripada sekitarnya: di atas laut yang hangat
kelembapan udara menambahkan gaya apung yang setara dengan beberapa kelvin, sama besarnya dengan kelebihan
suhu yang cukup besar.
\end{solution}

\begin{solution}{exo:b2:atmospheric-thermodynamics:6}
(a) $e_s$ = 6.11, 23.66, dan \qty{43.54}{hPa}; $r_s = 0.622\,e_s/(1000 - e_s)$ = 3.82, 15.07, dan
\qty{28.31}{g/kg}; $\Gamma_s$ = 6.48, 4.20, dan \qty{3.44}{K/km}. (b) Parsel kering mendingin
\qty{9.8}{K}, parsel jenuh \qty{4.2}{K}: parsel jenuh \qty{5.6}{K} lebih hangat.
(c) Makin ke atas udara makin dingin, $r_s$ makin kecil, tinggal sedikit uap yang dapat mengembun, dan
$\Gamma_s$ kembali mendekati $\Gamma_d$.
\end{solution}

\begin{solution}{exo:b2:atmospheric-thermodynamics:7}
(a) $e = e_s(\qty{291.65}{K}) = \qty{21.52}{hPa}$, $r = 0.622 \times 21.52/(971.1 - 21.52) =
\qty{14.1}{g/kg}$; $\theta = 292.15 \times (1000/971.1)^{0.2857} = \qty{294.6}{K}$; $\theta_e =
294.6 \exp[2.501 \times 10^6 \times 0.0141/(1004.5 \times 291.5)] = 294.6 \times 1.128 =
\qty{332.3}{K}$. (b) $e_s(\qty{264.25}{K}) = \qty{3.13}{hPa}$, $r_s = 0.622 \times 3.13/(500 -
3.13) = \qty{3.92}{g/kg}$; $\theta = 264.25 \times 2^{0.2857} = \qty{322.1}{K}$; parsel itu
jenuh, sehingga $T_{\text{LCL}} = T$ dan $\theta_e = 322.1 \exp[2.501 \times 10^6 \times
0.00392/(1004.5 \times 264.25)] = \qty{334.3}{K}$. (c) $\theta$ naik \qty{27.5}{K}, karena dipanaskan
oleh kalor laten dari \qty{10.2}{g/kg} uap yang mengembun; $\theta_e$ berubah
\qty{0.6}{\percent}, dalam batas ketelitian rumus hampirannya: besaran itu kekal.
\end{solution}

\begin{solution}{exo:b2:atmospheric-thermodynamics:8}
(a) $\text{CAPE} = 287.0 \times 4 \times \ln(700/300) = 287.0 \times 4 \times 0.847 =
\qty{973}{J/kg}$. (b) $\sqrt{2 \times 973} = \qty{44.1}{m/s}$. (c) Antara 1 dan
\qty{1500}{J/kg}: positif, tetapi tidak besar. (d) Pita dengan lebar horizontal tetap di antara dua
kurva; menurut \cref{prop:b2:atmospheric-thermodynamics:areas} luasnya sebanding dengan $\int
\Delta T\,\mathrm{d}\ln p = 4 \ln(700/300)$, berapa pun kemiringan isotermnya.
\end{solution}

\begin{solution}{exo:b2:atmospheric-thermodynamics:9}
(a) $D_c = \sqrt{3 \times 0.61 \times 80^3/2.1}\,\unit{nm} = \qty{668}{nm} =
\qty{0.668}{\micro m}$; $s_c = 2A/(3D_c) = 2 \times 2.1/(3 \times 668) = 0.00210 =
\qty{0.210}{\percent}$. (b) Selesaikan $s_c = \sqrt{4A^3/(27\kappa D_d^3)}$ untuk $D_d$: $D_d =
\bigl(4A^3/(27\kappa s_c^2)\bigr)^{1/3}$, yaitu \qty{63}{nm} pada \qty{0.3}{\percent} dan \qty{131}{nm} pada
\qty{0.1}{\percent}. (c) Setiap partikel antara 63 dan \qty{131}{nm} kini juga teraktivasi: lebih banyak
tetes, dan, untuk air cair yang sama banyak, tetes yang lebih kecil.
\end{solution}

\begin{solution}{exo:b2:atmospheric-thermodynamics:10}
(a) $w = \sqrt{2 \times 130} = \qty{16.1}{m/s}$, mustahil bagi udara yang diam dekat tanah.
(b) Suhu potensial puncak inversi adalah $293.65 \times (1000/889.1)^{0.2857} =
\qty{303.7}{K}$; parsel permukaan dengan $\theta$ ini memiliki $T = 303.7 \times (971.1/1000)^{0.2857}
= \qty{301.2}{K}$, yaitu \qty{28.0}{\celsius}. (c) $125 \times (28.0 - 18.5) = \qty{1190}{m}$ di atas
tanah, di atas puncak inversi yang lama. (d) Front, yang mengangkat udara di lerengnya; kenaikan lambat
udara di depan suatu gangguan di atas, yang mendinginkan tutup itu sendiri.
\end{solution}

\begin{solution}{exo:b2:atmospheric-thermodynamics:11}
(a) Bila $r_s \to 0$, pembilang mendekati $g$ dan penyebut mendekati $c_p$. (b) Tuliskan
$L^2r_s/(R_vT^2) = c_p \times \bigl(Lr_s/(R_dT)\bigr) \times LR_d/(R_vTc_p)$, dengan
$LR_d/(R_vTc_p) = 5.1$ pada \qty{303}{K}: bila $Lr_s/(R_dT) \gg 1$, juga $L^2r_s/(R_vT^2) \gg c_p$,
dan $\Gamma_s \to g\,(Lr_s/(R_dT))/(L^2r_s/(R_vT^2)) = gR_vT/(LR_d) = 9.81 \times 461.5 \times
303/(2.501 \times 10^6 \times 287.0) = \qty{1.91}{K/km}$. (c) Dengan $T_d = T$,
$\Gamma_{\text{embun}} = gR_vT/(LR_d)$, sama saja. Dalam limit itu kalor laten begitu besar sehingga
parsel hampir tidak dapat mengembunkan uapnya tanpa menghangat jauh di atas titik
kejenuhannya: parsel itu mempertahankan \omterm{def:b2:atmospheric-thermodynamics:moist-air}{rasio campurannya} dan mengikuti garis \omterm{def:b2:atmospheric-thermodynamics:moist-air}{rasio campuran} tetap, yang di sepanjangnya
suhu turun dengan laju $\Gamma_{\text{embun}}$. (d) Pada \qty{30}{\celsius} dan \qty{1000}{hPa},
$Lr_s/(R_dT) = 0.81$, tidak jauh lebih besar daripada 1, dan $L^2r_s/(R_vT^2) = 4180$, empat kali $c_p$:
$\Gamma_s = \qty{3.4}{K/km}$, masih jauh dari \qty{1.9}{K/km}. Limit itu tidak tercapai.
\end{solution}

\begin{solution}{exo:b2:atmospheric-thermodynamics:12}
(a) $2.868/2.599 = 1.104$ dan $1.257/1.032 = 1.218$: \qty{10.4}{\percent} dan
\qty{21.8}{\percent}. (b) Tetes teraktivasi pada \omterm{def:b2:atmospheric-thermodynamics:microphysics}{kelewatjenuhan} beberapa persepuluh persen,
puluhan kali lebih kecil. (c) Kristal tumbuh cepat melalui deposisi; kristal itu mengambil uap dari udara,
yang lalu turun di bawah kejenuhan terhadap tetes; tetes menguap dan uapnya pindah ke
es (proses Bergeron), sampai awan itu seluruhnya menjadi es. (d) Pada \qty{0}{\celsius} kedua
\omterm{def:b2:atmospheric-thermodynamics:moist-air}{tekanan uap} jenuh itu sama: keduanya bertemu di titik tripel, dan rasionya mendekati satu
di dekatnya.
\end{solution}

\begin{solution}{pb:b2:atmospheric-thermodynamics:1}
\textbf{1.} $e = e_s(\qty{291.65}{K}) = \qty{21.5}{hPa}$.
\textbf{2.} $r = 0.622 \times 21.5/(971.1 - 21.5) = \qty{14.1}{g/kg}$.
\textbf{3.} $q = 14.1/1.0141 = \qty{13.9}{g/kg}$.
\textbf{4.} $e_s(\qty{292.15}{K}) = \qty{22.2}{hPa}$; $21.5/22.2 = \qty{97}{\percent}$.
\textbf{5.} $0.61 \times 0.0141 \times 292.15 = \qty{2.5}{K}$.
\textbf{6.} $125 \times 0.5 = \qty{63}{m}$; $p = 971.1 \exp[-9.81 \times 62.5/(287.0 \times 294.7)]
= \qty{964}{hPa}$.
\textbf{7.} $\theta = 292.15 \times (1000/971.1)^{0.2857} = \qty{294.6}{K}$.
\textbf{8.} $16.9 - 17.3 = \qty{-0.4}{K}$: lebih dingin.
\textbf{9.} $(20.5 - 16.4)/(1102 - 924) = 4.1/178$: suhu naik \qty{23}{K/km}, yaitu
laju lapse \qty{-23}{K/km}. Sebuah inversi: stabil mutlak, sebuah tutup.
\textbf{10.} $15.4 - 20.5 = \qty{-5.1}{K}$.
\textbf{11.} $e_s(\qty{-1.8}{\celsius}) = \qty{5.4}{hPa}$ dan $e_s(\qty{18.6}{\celsius}) =
\qty{21.7}{hPa}$: \qty{25}{\percent}. Karena hangat dan kering, udara itu dapat berupa udara yang turun dan menghangat karena
pemampatan, atau udara yang dipanaskan di atas tanah kering di arah datangnya angin; udara itu bukan udara lembap di bawahnya.
\textbf{12.} $18.6 + 18.4 = \qty{37.0}{K}$ sepanjang $5725 - 1489 = \qty{4236}{m}$, \qty{8.73}{K/km}: di antara $\Gamma_s$ dan
$\Gamma_d$, tidak stabil bersyarat.
\textbf{13.} $-18.4 - (-8.9) = \qty{-9.5}{K}$: ketidakstabilan ekstrem ($-8$ sampai $-10$).
\textbf{14.} Dari tanah ke atas: 0, $-0.4$, $-0.2$, $-5.1$, $-4.9$, $+1.6$, $+9.5$, $+10.0$,
$+5.6$, $+1.5$, $-9.5$~\unit{K}.
\textbf{15.} Antara 850 ($-4.9$) dan \qty{700}{hPa} ($+1.6$), pada fraksi $4.9/6.5 = 0.754$:
$\ln p = \ln 850 + 0.754 \ln(700/850)$, $p = \qty{734}{hPa}$; $z = 1489 + 0.754 \times 1621 =
\qty{2711}{m}$, \qty{2.37}{km} di atas tanah.
\textbf{16.} Antara 250 ($+1.5$) dan \qty{200}{hPa} ($-9.5$), pada fraksi $1.5/11.0 =
0.136$: $p = \qty{242.5}{hPa}$ dan $z = 10560 + 0.136 \times 1446 = \qty{10757}{m}$, sekitar
\qty{200}{m} di bawah tropopause yang dilaporkan (\qty{10954}{m}).
\textbf{17.} Segmen-segmennya ($\Delta T$ rata-rata dikali $\ln$ rasio tekanan): $0.80 \times 0.0474
= 0.038$; $5.55 \times 0.3365 = 1.867$; $9.75 \times 0.2231 = 2.176$; $7.80 \times 0.2877 =
2.244$; $3.55 \times 0.1823 = 0.647$; $0.75 \times 0.0305 = 0.023$; jumlahnya \qty{6.995}{K}.
$\text{CAPE} = 287.0 \times 6.995 = \qty{2010}{J/kg}$, sampai tiga angka.
\textbf{18.} Antara 1500 dan \qty{2500}{J/kg}: besar.
\textbf{19.} $-0.2 \times 0.0486 - 0.3 \times 0.0189 - 2.65 \times 0.0207 - 5.0 \times 0.0450 -
2.45 \times 0.1467 = -0.655$~\unit{K}; $\text{CIN} = 287.0 \times (-0.655) = \qty{-188}{J/kg}$.
\textbf{20.} Tabel hanya memuat sebelas aras: di antara aras-aras itu aturan trapesium mengganti
kurva dengan segmen lurus dalam $\ln p$ dan melewatkan lengkungannya. Kisi \qty{1}{hPa} mengikuti kurva. \omterm{def:b2:atmospheric-thermodynamics:virtual}{Suhu virtual}
menambahkan gaya apung uap air parsel, yang lebih melimpah di dalam parsel jenuh daripada di
udara kering di sekitarnya: CAPE yang lebih besar, CIN yang lebih kecil.
\textbf{21.} $\sqrt{2 \times 2010} = \qty{63}{m/s}$.
\textbf{22.} $10757 - 2711 = \qty{8046}{m}$ dengan $63.4/2 = \qty{31.7}{m/s}$: \qty{254}{s}, sekitar empat
menit.
\textbf{23.} Berat air terkondensasi yang dibawa arus naik; percampuran dengan udara yang lebih kering
di sekitarnya (entrainment), yang mendinginkannya ketika tetes-tetes menguap.
\textbf{24.} Teori parsel memberikan kepada arus naik badai yang diberi makan oleh udara ini kecepatan sekitar
\textbf{\qty[text-series-to-math=true]{63}{m.s^{-1}}}, sebuah batas atas.
\end{solution}
